% ---------------------------------------------------------------------------
% macros.tex -- admonition boxes, page style, semantic shortcuts.
% ---------------------------------------------------------------------------

% --- semantic shortcuts -----------------------------------------------------
\newcommand{\SV}{SystemVerilog}
\newcommand{\VHDL}{VHDL}
\newcommand{\LRM}{LRM}
\newcommand{\ie}{i.e.\@\xspace}
\newcommand{\eg}{e.g.\@\xspace}
\usepackage{xspace}
\newcommand{\DUT}{DUT\xspace}
\newcommand{\RTL}{RTL\xspace}

% --- backtick helper, shared by every chapter -------------------------------
% Typesets a Verilog compiler directive: \tick{timescale 1ns/1ps}
\newcommand{\tick}[1]{\code{\textasciigrave#1}}
\newcommand{\chOneDir}[1]{\tick{#1}}
\newcommand{\chTwoTick}[1]{\tick{#1}}
\newcommand{\chFiveM}[1]{\tick{#1}}

% --- admonition boxes -------------------------------------------------------
% \begin{gotcha}{Short title}  ... \end{gotcha}
%     A trap that specifically catches VHDL designers.
\newtcolorbox{gotcha}[1]{%
  enhanced, breakable,
  colback=warnbg, colframe=warnframe, boxrule=0.9pt, arc=3pt,
  left=10pt,right=10pt,top=11pt,bottom=9pt,
  before skip=3.5mm, after skip=3.5mm,
  fonttitle=\bfseries\sffamily\small,
  coltitle=white, colbacktitle=warnframe,
  attach boxed title to top left={xshift=6mm,yshift=-3.2mm},
  boxed title style={arc=2pt,boxrule=0pt,left=6pt,right=6pt,
                     top=2.5pt,bottom=2.5pt},
  title={\raisebox{-0.5pt}{$\triangle$}\ Trap: #1}}

% \begin{tip}{Short title} ... \end{tip}
\newtcolorbox{tipbox}[1]{%
  enhanced, breakable,
  colback=tipbg, colframe=tipframe, boxrule=0.9pt, arc=3pt,
  left=10pt,right=10pt,top=11pt,bottom=9pt,
  before skip=3.5mm, after skip=3.5mm,
  fonttitle=\bfseries\sffamily\small,
  coltitle=white, colbacktitle=tipframe,
  attach boxed title to top left={xshift=6mm,yshift=-3.2mm},
  boxed title style={arc=2pt,boxrule=0pt,left=6pt,right=6pt,
                     top=2.5pt,bottom=2.5pt},
  title={Recipe: #1}}

% \begin{notebox}{Short title} ... \end{notebox}
\newtcolorbox{notebox}[1]{%
  enhanced, breakable,
  colback=notebg, colframe=noteframe, boxrule=0.9pt, arc=3pt,
  left=10pt,right=10pt,top=11pt,bottom=9pt,
  before skip=3.5mm, after skip=3.5mm,
  fonttitle=\bfseries\sffamily\small,
  coltitle=white, colbacktitle=noteframe,
  attach boxed title to top left={xshift=6mm,yshift=-3.2mm},
  boxed title style={arc=2pt,boxrule=0pt,left=6pt,right=6pt,
                     top=2.5pt,bottom=2.5pt},
  title={#1}}

% A short unnumbered "rule of thumb" strip.
\newtcolorbox{ruleofthumb}{%
  enhanced, breakable, sharp corners,
  colback=white, colframe=accent,
  boxrule=0pt, leftrule=3.5pt,
  left=11pt,right=8pt,top=7pt,bottom=7pt,
  before skip=3.5mm, after skip=3.5mm,
  fontupper=\itshape}

% Chapter-opening summary strip: \chapabstract{...}
\newcommand{\chapabstract}[1]{%
  \begin{tcolorbox}[enhanced,sharp corners,colback=codebg,colframe=codeframe,
                    boxrule=0.7pt,left=11pt,right=11pt,top=9pt,bottom=9pt,
                    before skip=2mm, after skip=7mm]
    \small\sffamily #1
  \end{tcolorbox}}

% --- figure placeholder -----------------------------------------------------
% \placeholder{height}{what the figure should show}
% Use it where a screenshot or an externally sourced drawing belongs.
\newcommand{\placeholder}[2]{%
  \begin{center}
  \begin{tikzpicture}
    \node[draw=codeframe,fill=codebg,line width=0.8pt,dashed,
          minimum width=0.82\linewidth,minimum height=#1,
          align=center,inner sep=8pt,rounded corners=2pt]
         {\sffamily\small\color{codenumber}
          [\,figure placeholder\,]\\[3pt]
          \parbox{0.72\linewidth}{\centering\footnotesize\itshape #2}};
  \end{tikzpicture}
  \end{center}}

% --- page style -------------------------------------------------------------
\pagestyle{fancy}
\fancyhf{}
\renewcommand{\headrulewidth}{0.4pt}
\fancyhead[LE]{\small\sffamily\nouppercase{\leftmark}}
\fancyhead[RO]{\small\sffamily\nouppercase{\rightmark}}
\fancyfoot[LE,RO]{\small\sffamily\thepage}
\fancypagestyle{plain}{\fancyhf{}\renewcommand{\headrulewidth}{0pt}%
  \fancyfoot[LE,RO]{\small\sffamily\thepage}}
\renewcommand{\chaptermark}[1]{\markboth{\chaptername\ \thechapter.\ #1}{}}
\renewcommand{\sectionmark}[1]{\markright{\thesection\ \ #1}}

% --- chapter headings -------------------------------------------------------
\titleformat{\chapter}[display]
  {\normalfont\sffamily\huge\bfseries\color{codekey}}
  {\raggedright\sffamily\normalsize\bfseries\color{codenumber}%
   \MakeUppercase{\chaptertitlename}\ \thechapter}
  {8pt}
  {\raggedright #1}
  [\vspace{-2pt}{\color{codeframe}\rule{\textwidth}{1.2pt}}]
\titlespacing*{\chapter}{0pt}{-10pt}{28pt}

\titleformat{\section}
  {\normalfont\sffamily\Large\bfseries\color{codekey}}{\thesection}{0.7em}{#1}
\titleformat{\subsection}
  {\normalfont\sffamily\large\bfseries}{\thesubsection}{0.6em}{#1}
\titleformat{\subsubsection}
  {\normalfont\sffamily\normalsize\bfseries\itshape}{\thesubsubsection}{0.6em}{#1}

% --- caption style ----------------------------------------------------------
\captionsetup{font=small,labelfont={bf,sf},labelsep=period,
              justification=raggedright,singlelinecheck=false}
\captionsetup[lstlisting]{font=small,labelfont={bf,sf},labelsep=period}
\renewcommand{\lstlistingname}{Listing}
\renewcommand{\lstlistlistingname}{List of Listings}

% --- misc -------------------------------------------------------------------
\setlength{\parskip}{0.35\baselineskip}
\setlength{\parindent}{0pt}
\raggedbottom
\setcounter{tocdepth}{1}
\setcounter{secnumdepth}{3}

% cleveref names
\crefname{lstlisting}{Listing}{Listings}
\Crefname{lstlisting}{Listing}{Listings}

% ---------------------------------------------------------------------------
% Cross-booklet references.  The SystemVerilog cookbook was split into four
% companion booklets (Design, Verification, Simulation, Example); a \cref
% that used to point at a chapter now living in a different booklet cannot
% resolve, so it is rewritten to name that booklet instead of a page number.
% ---------------------------------------------------------------------------
\newcommand{\refdesignbook}{the \emph{Design Cookbook}}
\newcommand{\refverificationbook}{the \emph{Verification Cookbook}}
\newcommand{\refsimulationbook}{the \emph{Simulation Cookbook}}
\newcommand{\refexamplebook}{the \emph{Example Cookbook}}
\newcommand{\Refdesignbook}{The \emph{Design Cookbook}}
\newcommand{\Refverificationbook}{The \emph{Verification Cookbook}}
\newcommand{\Refsimulationbook}{The \emph{Simulation Cookbook}}
\newcommand{\Refexamplebook}{The \emph{Example Cookbook}}

